\documentclass[10pt]{article}
\usepackage[utf8]{inputenc}
\usepackage[T1]{fontenc}
\usepackage{amsmath}
\usepackage{amsfonts}
\usepackage{amssymb}
\usepackage[version=4]{mhchem}
\usepackage{stmaryrd}
\usepackage{graphicx}
\usepackage[export]{adjustbox}
\graphicspath{ {./images/} }

\begin{document}
Adverse selection

\section*{Asymmetric information}
\begin{itemize}
  \item Asymmetry of information often breaks optimality properties of markets. Two main phenomena with many variations.
  \item Adverse selection (hidden information).
  \item Moral hazard (hidden action).
\end{itemize}

\section*{Adverse selection}
\begin{itemize}
  \item An agent's trading decisions depend on her private information. The informed agent's actions then reveal information. This affects, adversely, the uninformed side.
  \item Ex: Agents have private information about their innate health needs and then choose among insurance policies.
  \item Common models: signaling, screening, mechanism design in various models.
\end{itemize}

\section*{Moral hazard}
\begin{itemize}
  \item An agent's actions cannot be freely monitored, but affect other agent's utility.
  \item Example: The success of a firm depends on its manager's effort. Effort is unobservable by the owner. It may not be possible to induce the manager to provide the desired effort level.
  \item Common model: principal-agent.
\end{itemize}

\section*{Adverse selection}
\begin{itemize}
  \item Topics:
  \item Adverse selection and market inefficiency.
  \item Akerlof's "Market for lemmons".
  \item Second best outcome: Constrained optimality
  \item Signaling models
  \item Screening models
\end{itemize}

\section*{The market for lemons}
I wrote "The Market for Lemons," (a 13-page paper for which I was awarded the Prize in Economics) during my first year as assistant professor at Berkeley, in 1966-67. It concerns how horse traders respond to the natural question: "if he wants to sell that horse, do I really want to buy it?" Such questioning is fundamental to the market for horses and used cars, but it is also at least minimally present in every market transaction. - George Akerlof

George Akerlof, Writing the Market for Lemons: A Personal and Interpretive Essay.\\
link\_to\_essay

\section*{Market for lemons}
\begin{itemize}
  \item Objective: Asymmetry of information causes markets to fail. Often markets require regulation for efficiency.
  \item Market for used cars.
  \item Asymmetry of information: The seller has better information than the buyer about the "quality" of the car.
  \item See Akerlof (1970).
\end{itemize}

\section*{Example}
\begin{itemize}
  \item Depending on its quality a car is worth $x$ in $\{10,12,14,16,18,20\}$ where $x$ is in expressed in thousands.
  \item Seller knows quality $x$ of her car but buyer does not.
  \item Seller only sells if the price $p \geq x$.
  \item Buyer only buys if $p \leq E(x)$.
  \item All values of $x$ are equally likely.
  \item What is the market price $p$ and what type of cars will be traded?
\end{itemize}

Suppose $p>20$.

\begin{itemize}
  \item Sellers wish to sell since $(\forall x) p>20 \geq x$.
  \item Buyer does not buy:
\end{itemize}

$$
E(x)=\frac{1}{6} 10+\frac{1}{6} 12+\frac{1}{6} 14+\frac{1}{6} 16+\frac{1}{6} 18+\frac{1}{6} 20=15
$$

\begin{itemize}
  \item and $p>20>15=E(x)$.
\end{itemize}

Suppose $p=15$.

\begin{itemize}
  \item Sellers with $x \leq 15$ wish to sell. Sellers with $x>15$ do not.
  \item Buyers anticipate that sellers with $x>15$ do not sell.
\end{itemize}

$$
E(x)=\frac{1}{3} 10+\frac{1}{3} 12+\frac{1}{3} 14=12
$$

Since $p=15>12=E(x)$, buyers do not buy $(E(x)-p<0)$.

\section*{Example}
Suppose $p=10$.

\begin{itemize}
  \item Sellers of lemons $(x=10)$ wish to sell. Sellers with $x>10$ do not.
  \item Buyers anticipate that sellers with $x>10$ do not sell.
\end{itemize}

$$
E(x)=\frac{1}{1} 10=10
$$

\begin{itemize}
  \item Buyers buy lemons; $p=10=E(x)$, knowing what they buy.
\end{itemize}

\section*{Model}
\begin{itemize}
  \item The seller observes her type $x \in\left[x^{\prime}, x^{\prime \prime}\right] ; x$ is the value to the seller.
  \item The yype $x$ distributed according to density $f(x)$.
  \item The buyer does not observe $x$.
  \item The buyer is willing to pay $v(x)$ for car $x$; buyer is expected utility maximizer.
  \item Positive gains from trade: $\left(\forall x \in\left(x^{\prime}, x^{\prime \prime}\right]\right) v(x)>x$ and $v\left(x^{\prime}\right)=x^{\prime}$.
  \item The functions $f$ and $v$ and the structure of the model are common knowledge.\\
(Both buyer and seller believe $x$ is distributed according to $f$, common prior assumption.)
\end{itemize}

\section*{Benchmark: Complete information}
\begin{itemize}
  \item $x$ is publicly observed before any transaction is arranged.
  \item Market outcome:
  \item Seller sells only if $p \geq x$.
  \item Buyer buys only if $v(x) \geq p$.
  \item Provided $v(x) \geq p \geq x$, trade takes place.
  \item Gains from trade are realized and divided among participants.
\end{itemize}

\section*{Incomplete information}
The buyer's expected value of a car depends on the buyer's beliefs about the cars in the market.

If types in market are $\left[x^{\prime}, x\right] \subseteq\left[x^{\prime}, x^{\prime \prime}\right]$, the buyer's expected value of a car is

$$
B(x)=\frac{\int_{x^{\prime}}^{x} v(z) f(z) d z}{\int_{x^{\prime}}^{x} f(z) d z}
$$

\begin{itemize}
  \item Set price $p$. Then $x>p \Longrightarrow[x$ is not in the market $]$.
  \item The buyer's expected value of the car when price is $p$
\end{itemize}

$$
B(p)=\frac{\int_{x^{\prime}}^{p} v(z) f(z) d z}{\int_{x^{\prime}}^{p} f(z) d z}
$$

\section*{Graph, lemons example}
\begin{center}
\includegraphics[max width=\textwidth, alt={}]{12d3571a-bf70-4859-959b-c133841dd5ee-17_1460_2578_347_155}
\end{center}

\section*{Graph, lemons example 2}
\begin{center}
\includegraphics[max width=\textwidth, alt={}]{12d3571a-bf70-4859-959b-c133841dd5ee-18_1457_2573_350_157}
\end{center}

\section*{Graph, Not lemons}
\begin{center}
\includegraphics[max width=\textwidth, alt={}]{12d3571a-bf70-4859-959b-c133841dd5ee-19_1408_2573_347_157}
\end{center}

\section*{Other possibilities}
\begin{center}
\includegraphics[max width=\textwidth, alt={}]{12d3571a-bf70-4859-959b-c133841dd5ee-20_1459_1363_350_162}
\end{center}

\section*{Constrained Pareto optimal solution}
\begin{center}
\includegraphics[max width=\textwidth, alt={}]{12d3571a-bf70-4859-959b-c133841dd5ee-21_1521_1773_345_159}
\end{center}

\section*{Responses to adverse selection}
-What can the seller do?\\
The seller, the informed party, can try to signal her information. (Signaling model)

\begin{itemize}
  \item What can the buyer, the uninformed party, do?
\end{itemize}

Could the buyer buy the information he is missing from the seller? How?

\section*{Signaling and screening}
\begin{itemize}
  \item Adverse selection leads to inefficiencies. Sometimes these inefficiencies originate on the inability of the uninformed party to separate the different types.\\
Signaling or screening practices.
  \item Signaling. The good type tries to separate from the bad type using signals.
  \item Screening. The uninformed party proposes contracts, such as multiple product lines, to screen consumers of different types.
\end{itemize}

Ex. Third degree price discrimination, insurance, etc.

\section*{References}
Akerlof, George A. 1970. "The Market for "Lemons": Uncertainty and the Market Mechanism." Quarterly Journal of Economics 84: 488-500.


\end{document}